%======================================================================
%  Guide: rg_session_fenics.py (FEniCS / DOLFINx)
%======================================================================
\sessiontitle{Guide 2 -- FEniCS: \texttt{rg\_session\_fenics.py}}
  {Blocks, figures and how to run them. Built on
   \texttt{reciprocity\_gap\_crack.py} (DOLFINx + gmsh).}

\untested{Status: not run with DOLFINx yet. Blocks 1 and 4--9 were run
(gmsh mesh, then the same arrays assembled with scikit-fem: this is the
Python file of Guide 3); Blocks 2--3 use the DOLFINx calls of
reciprocity\_gap\_crack.py. The values of Section 5 of the common part are
the references.}

\guidesec{1\quad Running the file}
\begin{enumerate}[itemsep=1pt]
\item Install DOLFINx 0.8--0.10 with gmsh (conda is the simplest route):
\begin{lstlisting}[style=shstyle]
conda create -n fenicsx -c conda-forge fenics-dolfinx mpich gmsh python-gmsh scipy matplotlib
conda activate fenicsx
\end{lstlisting}
\item Run it in serial, with the noise level and the folder of the figures
as optional arguments:
\begin{lstlisting}[style=shstyle]
python3 rg_session_fenics.py              # eps = 0, figures in .
python3 rg_session_fenics.py 0.01 figs    # eps = 0.01, figures in figs/
RG_NELEM=60 python3 rg_session_fenics.py  # finer mesh
\end{lstlisting}
\item The figures are written to \texttt{fig1.pdf} \dots\ \texttt{fig9.pdf}; the
messages \texttt{FIG k : ...} are printed when each is saved. The run takes a
few seconds.
\end{enumerate}

\guidesec{2\quad Parameters (Block 0)}
\begin{center}\small
\begin{tabular}{>{\ttfamily}l l p{8cm}}
\toprule
\rmfamily name & default & meaning\\
\midrule
ldom, hdom & 1.5, 2.0 & width and height of the rectangle\\
nelem & 30 & mesh size $\ell/\texttt{nelem}$, $\ell_\Gamma/(\texttt{nelem}/2)$ on the crack (\texttt{RG\_NELEM})\\
A, B & crack tips & the true crack (mesh and errors only)\\
eps & 0 & relative noise on the measured values (first argument)\\
seed & 1 & seed of the noise: change it to draw again\\
nxi, xiest & 8, 2.0 & Fourier fields $\xi=1\dots\texttt{nxi}$; $\xi$ of the length estimate\\
nsine & 5 & number of sine--sinh terms\\
nbeta & 12 & turning potential: angles $0,15,\dots,165$ degrees\\
outdir & . & folder of the figures (second argument)\\
\bottomrule
\end{tabular}
\end{center}

\guidesec{3\quad The blocks}
\block{Block 1 -- Geometry and mesh.} As in \texttt{replica.py}: gmsh
rectangle, the crack line embedded (\texttt{mesh.embed}), size
\texttt{Mesh.MeshSizeMax} $=\ell/30$ and $\ell_\Gamma/15$ at the tips. The gmsh
plugin \texttt{Crack} then duplicates the nodes of the crack (not the tips).
\texttt{gmshio.model\_to\_mesh} gives the DOLFINx mesh (the module is
\texttt{dolfinx.io.gmsh} from 0.10).

\block{Block 2 -- Model.} \texttt{functionspace(msh, ("Lagrange", 2))}: P2,
the TRI6 of Cast3M. The matrix $\mat K$ is assembled \emph{without} boundary
conditions and converted to SciPy (\texttt{fem.assemble\_matrix(a).to\_scipy()}):
$q=\mat Ku$ are then the nodal fluxes, as \texttt{KK * u} in Cast3M. The node
coordinates are \texttt{V.tabulate\_dof\_coordinates()}; \texttt{upnode} marks
the nodes of the cells on the side of $\vn$, which tells the lips apart.
FIG~1 uses the vertices of the cells (P1 space).

\block{Block 3 -- Procedures.} \texttt{rgap}, \texttt{measure}, \texttt{opening}
work on NumPy arrays and are the same as in Python. \texttt{dirichlet(mask,
values)} solves with \texttt{LinearProblem} and \texttt{dirichletbc} on the
nodes of \texttt{mask} (direct LU solver); the other boundary nodes and the
lips get the natural zero flux.

\block{Blocks 4--9.} The same text as in the Python file (Guide~3): from Block 4
on, the identification needs only the boundary values, the matrix $\mat K$
and the node coordinates. E1 is \texttt{dirichlet(bot | top, y)}; the turning
potential uses \texttt{dirichlet(outer, X)} and \texttt{dirichlet(outer, Y)}.
$R$ is decomposed with \texttt{numpy.linalg.svd}.

\guidesec{4\quad The figures}
\begin{center}\small
\begin{tabular}{>{\sffamily}l l l}
\toprule
figure & block & matplotlib\\
\midrule
FIG 1 & 2 & \texttt{triplot} of the vertices, boundary and crack\\
FIG 2 & 4 & \texttt{tricontourf} of $u_1$ at the vertices\\
FIG 3 & 4 & opening on the lips and infinite-body curve\\
FIG 4, 5 & 7 & identified line and segment; semi-ellipses\\
FIG 6, 7 & 8 & $C/m_0$, $S/m_0$ against $\xi$; sine--sinh sums\\
FIG 8, 9 & 9 & $|m_0|$ and the columns of $R$ against $\beta$\\
\bottomrule
\end{tabular}
\end{center}
The mesh is the gmsh mesh of the Python file: the figures are those of the
common part (Section 4).

\guidesec{5\quad The printed messages}
One line per block, \texttt{BLOCK k ...}, with the values of Section 5 of the
common part (same mesh, same element: the same numbers up to round-off).
Signs: $\vn$ is oriented so that $m_0>0$.

\guidesec{6\quad FEniCS pitfalls}
\begin{itemize}[itemsep=1pt]
\item \textbf{The lips are boundary facets.} After the \texttt{Crack} plugin,
\texttt{locate\_entities\_boundary} also returns the facets of the lips. Here
the outer boundary is selected by coordinates (\texttt{outer}), never by
``boundary facet''.
\item \textbf{Reactions.} \texttt{LinearProblem} with Dirichlet conditions
modifies the rows of the constrained nodes; the fluxes $\mat Ku$ need the
matrix assembled without conditions (Block 2).
\item \textbf{Serial only.} With MPI, the arrays are distributed and the
indices of \texttt{tabulate\_dof\_coordinates} are local; run with one
process.
\item \textbf{Versions.} \texttt{dolfinx.io.gmshio} became \texttt{dolfinx.io.gmsh}
in 0.10, and \texttt{LinearProblem} needs \texttt{petsc\_options\_prefix} there;
the file tries both.
\item \textbf{Mesh regularity.} The Fourier and sine--sinh fields are not
polynomials: their gaps carry a discretization error that depends on the size
\emph{and} the regularity of the mesh (sheet, question f.4).
\end{itemize}

\guidesec{7\quad Listing of \texttt{rg\_session\_fenics.py}}
\lstinputlisting[style=pystyle,lastline=231]{fenics/ch3/rg_session_fenics.py}
\noindent\emph{Blocks 4--9: the same text as in the Python file (Guide~3,
Section~7).}
